\documentclass[12pt,a4paper]{article}

\usepackage[T2A]{fontenc}
\usepackage[utf8]{inputenc}
\usepackage[russian]{babel}
\usepackage{geometry}
\usepackage{setspace}
\usepackage{indentfirst}
\usepackage{hyperref}
\usepackage{longtable}
\usepackage{array}

\geometry{top=2cm,bottom=1.5cm,left=3cm,right=2cm}
\onehalfspacing

\hypersetup{
  colorlinks=true,
  linkcolor=black,
  urlcolor=blue,
  pdftitle={Пояснительная записка: Эмулятор CPU 6502},
  pdfauthor={Ступин Александр Юрьевич}
}

\begin{document}

\begin{titlepage}
\begin{center}
\small
\textbf{Муниципальное бюджетное общеобразовательное учреждение\\
«Гимназия №2 «Квантор»}}\\[0.7cm]

\textbf{Индивидуальный проект}\\
по предмету: \textit{Информатика}\\
тип проекта: \textit{учебный, практико-ориентированный}\\[1.2cm]

{\LARGE \textbf{Разработка учебного веб-приложения\\
«Эмулятор CPU 6502 с ассемблером, консолью и пошаговой трассировкой»}}\\[1.5cm]

\begin{flushright}
\begin{tabular}{l}
Выполнил: Ступин Александр Юрьевич \\
Класс: 10 «В» \\
Руководитель проекта: Пименова Анна Николаевна
\end{tabular}
\end{flushright}

\vfill
\textbf{Коломна --- 2026}
\end{center}
\end{titlepage}

\tableofcontents
\newpage

\section{Введение}

\subsection{Формулирование проблемы}
В школьном курсе информатики изучение архитектуры процессора чаще всего ограничивается теорией. Учащимся сложно увидеть на практике, как выполняются инструкции ассемблера, как изменяются регистры, флаги и память после каждого шага программы. Существующие инструменты часто либо слишком сложны для учебных целей, либо не дают удобной визуализации для анализа.

Современный ученик привык к быстрой обратной связи от приложений и игр: любое действие сопровождается визуальным ответом. В то же время классическое объяснение темы «как работает процессор» на доске или по учебнику не даёт того же уровня вовлечённости. В результате формируется иллюзия, что «настоящее программирование» --- это только языки высокого уровня, а низкий уровень остаётся чем-то далёким и необязательным. На практике же понимание регистров, памяти и условных переходов полезно не только будущим инженерам, но и всем, кто хочет глубже понимать работу компьютера в целом.

Отдельную трудность представляет подбор инструмента: профессиональные среды разработки перегружены функциями, требуют длительной настройки и отвлекают от сути --- от пошагового разбора программы. Учебные эмуляторы в браузере зачастую либо упрощены до минимума, либо наоборот ориентированы на энтузиастов ретро-техники и не адаптированы под формат урока в общеобразовательной школе. Таким образом, проблема заключается в отсутствии сбалансированного решения: достаточно наглядного, но при этом доступного с первых минут работы.

\subsection{Гипотеза проекта}
Если разработать учебное веб-приложение, объединяющее эмулятор CPU 6502, ассемблер и пошаговую визуализацию состояния процессора, то понимание учащимися принципов низкоуровневого программирования и архитектуры процессора станет более глубоким и осознанным.

Гипотеза опирается на общепедагогическую идею: человек лучше запоминает то, что не только прочитал или услышал, но и сам «прожил» в виде последовательности наблюдаемых состояний. Когда ученик видит, что после конкретной инструкции изменился флаг нуля или содержимое ячейки памяти, абстрактное правило превращается в конкретный опыт. Веб-формат при этом снимает барьер установки ПО: достаточно открыть страницу в браузере на любом компьютере класса или дома, что особенно важно для равного доступа к образовательным ресурсам.

\subsection{Актуальность проекта}
Актуальность проекта обусловлена потребностью в доступном интерактивном инструменте для изучения темы «Архитектура ЭВМ». Наглядный пошаговый разбор исполнения программ позволяет сократить разрыв между теорией и практикой, повышает интерес учащихся к системному программированию и делает изучение ассемблера более осмысленным.

В контексте формирования функциональной грамотности и цифровой культуры особенно ценны задания, где ученик не просто получает готовый ответ, а строит цепочку рассуждений: почему программа пошла по одной ветке условного перехода, а не по другой; как связаны арифметические операции и флаги переноса. Такие навыки переносимы на более широкий круг задач --- от отладки алгоритмов до понимания ограничений вычислительной техники. Кроме того, тема архитектуры ЭВМ исторически и методически связана с развитием логического мышления, что делает её содержательно значимой даже вне узкой специализации на IT.

Актуальность подчёркивается и тем, что 6502 --- это «классика», наглядно демонстрирующая принципы работы процессора без избыточной сложности современных многоядерных систем. Для школьного уровня это компромисс между реалистичностью модели и возможностью уложить объяснение в ограниченное время урока.

\subsection{Объект исследования}
\textbf{Объект исследования} --- процесс обучения архитектуре ЭВМ и основам ассемблерного программирования в школьном курсе информатики.

Под объектом здесь понимается не отдельный урок и не один раздел учебника, а целостный образовательный процесс, в котором учащийся постепенно формирует представление о том, как машина исполняет программу. Важны и мотивация, и темп, и возможность повторить материал дома без специального оборудования. Именно поэтому объект исследования описывается широко: речь идёт о сопровождении учебной темы цифровым инструментом, который можно встроить в разные форматы работы --- фронтальное объяснение у учителя, самостоятельная работа, проектная деятельность.

\subsection{Предмет исследования}
\textbf{Предмет исследования} --- методы разработки и педагогическая эффективность интерактивного программного продукта для визуализации исполнения инструкций CPU.

В рамках предмета исследования рассматриваются вопросы выбора технологий, разделения ответственности между клиентской и серверной частью, организации данных (трасса, состояние, выводы) и того, насколько интерфейс поддерживает типичный учебный сценарий: написал код --- собрал --- запустил --- объяснил по шагам. Педагогическая эффективность здесь оценивается не формальным тестом, а логикой проектирования: насколько продукт снижает когнитивную нагрузку и направляет внимание на ключевые элементы модели процессора.

\subsection{Цели проекта}
Разработать и апробировать учебное веб-приложение «Эмулятор CPU 6502», позволяющее вводить программу на ассемблере, запускать её и пошагово анализировать исполнение команд по состоянию регистров, флагов и памяти.

Формулировка цели намеренно включает два аспекта: \textit{разработать} (то есть создать работоспособный программный продукт) и \textit{апробировать} (то есть проверить применимость в учебной логике --- на примерах программ, в типовых сценариях запуска, при работе с консолью и трассой). Такое разделение важно для проектной деятельности: итог оценивается не только «запускается ли программа», но и «удобно ли ею пользоваться в роли ученика или учителя».

\subsection{Задачи проекта}
Задачи проекта выстроены в логической последовательности от теории к практике и от «ядра» к «оболочке». Такой порядок отражает реальный ход инженерной работы: сначала необходимо понять предметную область, затем выбрать инструменты, спроектировать взаимодействие компонентов, реализовать вычислительную часть, после чего --- интерфейс и интеграцию, и лишь затем перейти к проверке качества и осмыслению результата.

\begin{enumerate}
  \item Изучить архитектуру и базовые принципы работы процессора 6502.
  \item Проанализировать подходящие технологии реализации клиентской и серверной частей.
  \item Спроектировать архитектуру приложения (frontend + backend + протокол обмена).
  \item Реализовать эмулятор CPU и ассемблер в серверной части на Python.
  \item Реализовать браузерный интерфейс на React + TypeScript.
  \item Настроить API-взаимодействие между frontend и backend.
  \item Провести тестирование основных пользовательских сценариев.
  \item Оценить образовательную ценность и сформулировать перспективы развития.
\end{enumerate}

Отдельно стоит подчеркнуть, что задачи не сводятся к «написать код ради кода»: тестирование и рефлексия по итогам --- обязательные элементы проектной культуры. Они позволяют увидеть слабые места интерфейса (например, неочевидность сообщений об ошибках) и зафиксировать направления улучшения без смешения их с основной разработкой.

\subsection{Целевая группа}
Целевая группа проекта:
\begin{itemize}
  \item учащиеся 9--11 классов, изучающие информатику;
  \item учителя информатики, использующие демонстрационные материалы на уроках;
  \item начинающие разработчики, интересующиеся низкоуровневым программированием.
\end{itemize}

Для учащихся продукт выступает в роли тренажёра: можно экспериментировать с короткими программами, не рискуя повредить реальную систему, и при этом получать объяснимую трассу. Для учителя важна демонстрационная ясность --- возможность на большом экране показать, как одна инструкция сменяет другую, и ответить на вопросы «что будет, если поменять условие?». Для начинающих разработчиков проект полезен как пример связки современного frontend-стека с «классической» вычислительной задачей, что хорошо ложится в портфолио и расширяет кругозор за пределы типовых веб-форм.

При этом целевая группа не ограничивается перечислением: любой пользователь, которому интересна тема устройства процессора, может воспользоваться приложением в ознакомительных целях. Такая универсальность --- следствие выбора веб-платформы и отсутствия жёсткой привязки к конкретному учебнику или параграфу, хотя методически продукт ориентирован именно на школьный контекст.

С практической точки зрения полезно зафиксировать, что проект рассматривается как \textit{учебный программный продукт}, а не как коммерческий сервис. Это означает: приоритет отдаётся понятности, устойчивости сценариев демонстрации и воспроизводимости запуска на типичном школьном компьютере, а не максимальному количеству «маркетинговых» функций. Такая рамка помогает автору сосредоточиться на качестве ключевого пользовательского пути и не распыляться на второстепенные улучшения, которые не усиливают образовательный эффект.

Отдельно стоит упомянуть роль саморефлексии в проектной деятельности. По ходу работы приходилось неоднократно возвращаться к формулировке цели и проверять, соответствует ли очередное изменение кода изначальному замыслу. Иногда это приводило к отказу от «красивой, но лишней» идеи в пользу более простого и надёжного решения. Для пояснительной записки такой опыт важен тем, что он показывает зрелость подхода: проект --- это не только код, но и умение планировать, корректировать курс и аргументировать принятые решения.

\section{Описание проекта}

\subsection{Описание методов исследования}
В процессе работы использованы следующие методы:
\begin{itemize}
  \item \textbf{Анализ литературы и документации}: изучение источников по архитектуре 6502, Python, React, TypeScript и веб-разработке.
  \item \textbf{Проектный метод}: постановка цели, декомпозиция задач, последовательная реализация модулей.
  \item \textbf{Модульный подход}: раздельная разработка эмулятора, API и интерфейса.
  \item \textbf{Тестирование}: проверка корректности сборки и исполнения программ, анализа трассы и работы консоли ввода/вывода.
\end{itemize}

Метод анализа литературы применялся на протяжении всего цикла: на старте --- для выбора адекватной модели процессора и понимания набора инструкций, в середине --- для уточнения деталей реализации и согласования терминологии с общепринятыми обозначениями, на финишной прямой --- для корректного оформления записки и ссылок на авторитетные источники. Такой «сквозной» характер анализа снижает риск ошибок из-за недопонимания и помогает удерживать единый стиль изложения как в коде, так и в документации.

Проектный метод проявился в том, что работа не велась хаотично: сначала фиксировались ожидаемые пользовательские действия (что должен уметь делать ученик), затем под них подбиралась архитектура. Это типичный приём инженерии требований в упрощённом виде: даже в школьном проекте полезно явно отделить «хотелки» от минимально необходимого результата. Модульный подход, в свою очередь, облегчал отладку: ошибки сборки текста программы и ошибки отрисовки интерфейса рассматривались в разных слоях системы, что ускоряет поиск причины.

Тестирование в учебном проекте выполняло роль не только контроля качества, но и обучающего инструмента для самого автора: приходилось составлять короткие эталонные программы, проверять граничные случаи (пустой ввод, некорректная строка, останов по лимиту шагов) и наблюдать, насколько понятны сообщения об ошибках со стороны пользователя, который видит интерфейс впервые.

\subsection{Анализ используемых источников}
Для реализации проекта были использованы официальные технические источники и учебные материалы. Их можно разделить на группы:
\begin{enumerate}
  \item \textbf{Документация по инструментам разработки}: Python Documentation, React Documentation, TypeScript Documentation, Vite Guide.
  \item \textbf{Источники по архитектуре 6502}: Programming manuals и справочники по инструкциям.
  \item \textbf{Методические и образовательные материалы}: источники по организации проектной деятельности и формированию структуры пояснительной записки.
\end{enumerate}

Анализ показал, что для учебного проекта оптимальна full-stack архитектура: Python для вычислительной логики эмулятора и React + TypeScript для наглядного интерфейса.

Отдельного внимания заслуживает группа источников по Python. Язык выбран не только как популярный инструмент, но и как средство быстрой проверки идей: читаемый синтаксис облегчает сопровождение кода, а обширная стандартная библиотека и экосистема упрощают организацию серверной части. Документация React и TypeScript помогла выстроить интерфейс декларативно и заранее продумать типы данных, которые приходят с сервера (например, структура трассы и состояния регистров). Vite в данном контексте --- это скорее инфраструктурный выбор ускоренной разработки, но он также влияет на удобство демонстрации проекта: быстрый перезапуск локального сервера экономит время при итерациях интерфейса.

Источники по 6502 выполняют двойную функцию: с одной стороны, задают «правила игры» для эмулятора (что считать корректным поведением инструкции), с другой --- помогают корректно сформулировать справочные пояснения в интерфейсе. Наконец, методические материалы по проектной деятельности полезны для соблюдения логики изложения записки и для того, чтобы защита проекта выглядела цельной: не только «что сделано», но и «зачем», «как проверено», «что дальше».

\subsection{Описание хода работы над проектом}
Работа над проектом выполнялась по этапам:
\begin{enumerate}
  \item Определение темы, проблемы, цели и требований к итоговому продукту.
  \item Изучение предметной области: архитектуры 6502, команд, регистров, флагов, принципов исполнения.
  \item Проектирование структуры приложения: backend (эмулятор, ассемблер, API) и frontend (редактор, запуск, визуализация).
  \item Реализация backend на Python и настройка HTTP JSON API для сборки и запуска программ.
  \item Реализация frontend на React + TypeScript: редактор кода, вывод байтов программы, консоль, step inspector, ручной ввод при команде CTA.
  \item Интеграция частей приложения и отладка обмена данными через `/api`.
  \item Тестирование на типовых программах, корректировка интерфейса и улучшение UX.
  \item Финальная подготовка пояснительной записки и материалов к защите.
\end{enumerate}

На этапе постановки требований особое значение имело согласование ожиданий: продукт должен оставаться учебным, то есть не превращаться в профессиональный отладчик со сложной конфигурацией. Поэтому интерфейс сознательно ориентировался на понятные блоки: редактор, запуск, консоль, трассировка. На этапе изучения предметной области пришлось систематизировать знания о регистрах и флагах, поскольку именно они чаще всего вызывают вопросы у начинающих: что означает очередное значение флага, почему изменился счётчик команд, как связаны переходы и результат сравнения.

Проектирование структуры включало решение о разделении ответственности: вычисления и разбор текста программы выполняются на сервере, а визуализация --- в браузере. Такой подход типичен для веб-систем и упрощает сопровождение: логику эмуляции проще тестировать отдельно от графического оформления. Реализация API потребовала продумать форматы запросов и ответов, а также сценарий интерактивного ввода, когда программа останавливается в ожидании данных и продолжается после отправки значения пользователем.

Интеграция frontend и backend сопровождалась типичными для веб-разработки вопросами: корректность URL, проксирование запросов в режиме разработки, обработка ошибок сети и сервера. Эти моменты не являются «учебной теорией» в узком смысле, но они важны для реальной работоспособности продукта и для того, чтобы проект можно было уверенно демонстрировать на защите. UX-доработки на финальных этапах касались читаемости подсказок, логичности расположения элементов и визуальной согласованности светлой и тёмной темы.

Итоговый продукт:
Проектный продукт --- учебное веб-приложение «Эмулятор CPU 6502», включающее:
\begin{itemize}
  \item текстовый редактор ассемблерного кода;
  \item модуль сборки/запуска программы;
  \item консоль ввода и вывода значений;
  \item пошаговый просмотр исполнения (STEP) с навигацией стрелками;
  \item отображение регистров A, X, Y, PC, OP и блока MEMORY;
  \item светлую и тёмную темы интерфейса;
  \item справочный раздел по инструкциям.
\end{itemize}

Ниже приведено более развёрнутое описание содержательной и технической стороны продукта, поскольку именно сочетание «понятного интерфейса» и «корректной модели исполнения» определяет ценность учебного приложения.

\textbf{Краткий историко-методический контекст.} Процессор 6502 занимает заметное место в истории вычислительной техники: он использовался в ряде массовых систем, и его архитектура до сих пор служит удачным введением в тему машинного уровня. Для школьного проекта это удобно тем, что можно опереться на большое количество справочных материалов и при этом не перегружать ученика чрезмерным количеством режимов адресации и периферийных особенностей, характерных для более молодых платформ.

\textbf{Пользовательский сценарий.} Типичная работа с приложением выглядит следующим образом. Пользователь вводит текст программы в редакторе, при необходимости сверяясь со встроенным справочником по инструкциям. Затем выполняется сборка: система либо формирует последовательность байт, либо возвращает сообщение об ошибке, которое важно трактовать как обучающую обратную связь, а не как «наказание». После успешной сборки запускается исполнение с ограничением по числу шагов, что защищает от бесконечных циклов в учебной среде. Если программа требует ввода, пользователь получает понятное состояние ожидания и вводит значение в консоли. Параллельно доступен просмотр трассировки: можно шаг за шагом наблюдать, как меняются регистры и какие ячейки памяти затрагиваются, что особенно полезно при разборе циклов и условных переходов.

\textbf{Архитектура и обмен данными.} Клиентская часть реализована как одностраничное приложение и отвечает за отображение состояния, ввод текста и отправку запросов к API. Серверная часть на Python выполняет разбор исходного текста, эмуляцию и формирование структурированного ответа, который затем отображается в интерфейсе. Такое разделение соответствует распространённой практике веб-разработки и упрощает дальнейшее развитие: например, можно добавить сохранение сессий или авторизацию, не переписывая ядро эмулятора.

\textbf{Организация работы и техника безопасности.} Поскольку проект выполнялся преимущественно за персональным компьютером, отдельное внимание уделялось режиму труда: перерывы, освещение, положение монитора. С точки зрения информационной безопасности учебный продукт не предполагает хранения персональных данных пользователей в составе минимальной версии, однако при публикации в интернете необходимо учитывать общие рекомендации: обновления зависимостей, корректная настройка прокси-сервера и ограничение поверхности атаки на стороне backend.

\textbf{Ограничения и честность постановки.} Любой эмулятор в учебном контексте является моделью, а не абсолютно полной копией железа. Поэтому в записке важно подчеркнуть: цель --- не заменить промышленные инструменты разработчика встроенных систем, а дать наглядный тренажёр для освоения базовых понятий. Именно такая рамка делает проект методически честным и снимает завышенные ожидания относительно полноты аппаратной совместимости.

\section{Заключение}
\subsection{Вывод. Достигнуты ли цели?}
По итогам выполнения работы цель проекта достигнута: разработано учебное веб-приложение «Эмулятор CPU 6502», в котором реализованы:
\begin{itemize}
  \item ввод программы на ассемблере;
  \item сборка и запуск программы;
  \item интерактивный ввод данных во время исполнения (CTA);
  \item вывод результатов выполнения;
  \item пошаговая трассировка с отображением регистров, флагов и памяти.
\end{itemize}

Поставленные задачи выполнены в полном объёме: изучена предметная область, реализована архитектура frontend + backend, проведено тестирование основных сценариев. Гипотеза проекта подтверждена: интерактивный формат действительно улучшает понимание принципов работы процессора и низкоуровневых программ.

В заключительной части уместно подвести итог не только по функциональному списку, но и по смыслу проделанной работы для автора. Разработка учебного приложения потребовала синхронизации двух разных плоскостей знаний: «железная логика» эмуляции и «человеческая логика» интерфейса. Именно на стыке этих плоскостей возникает ценность продукта: если интерфейс красивый, но модель неверна, обучение становится вредным; если модель верна, но интерфейс запутан, обучение не состоится из-за фрустрации пользователя. В рамках проекта удалось выдержать баланс на уровне, достаточном для демонстрации и учебного применения.

Отдельно отметим методический аспект: пошаговая трассировка превращает абстрактное «исполнение программы» в последовательность наблюдаемых состояний. Для ученика это снижает порог входа: не нужно сразу держать в голове всю программу целиком, можно двигаться маленькими шагами и сверять гипотезы с фактическими изменениями регистров. Для учителя это удобный инструмент сопровождения устного объяснения: можно остановиться на спорном месте и показать, что именно произошло на конкретном шаге.

Наконец, достижение цели подтверждается воспроизводимостью сценария: приложение может быть запущено локально, продемонстрировано на уроке и использовано для самостоятельной практики. Это важный практический критерий для школьного проекта --- результат не должен существовать только «в теории» на слайдах, он должен быть доступен как работающий артефакт.

\subsection{Ожидаемые результаты и перспективы развития}
Ожидаемый образовательный эффект проекта:
\begin{itemize}
  \item повышение качества усвоения темы «Архитектура ЭВМ»;
  \item развитие практических навыков работы с ассемблером;
  \item повышение мотивации учащихся к изучению системного программирования.
\end{itemize}

Ожидаемые результаты следует понимать не только как «оценки станут выше» (хотя улучшение успеваемости возможно), но и как изменение качества понимания: ученик начинает связывать причину и следствие между кодом и состоянием машины, а учитель получает инструмент для наглядной аргументации. В долгосрочной перспективе такой навык поддерживает интерес к информатике как к науке о системах, а не только как к набору прикладных программ.

Перспективы развития проекта:
\begin{itemize}
  \item расширение набора инструкций и режимов адресации;
  \item добавление breakpoints и расширенного режима отладки;
  \item визуализация карты памяти в hex-формате;
  \item внедрение набора учебных заданий с автоматической проверкой;
  \item публикация стабильной онлайн-версии для использования на уроках.
\end{itemize}

Среди перспектив можно также назвать развитие методического сопровождения: короткие рабочие листы с заданиями, скринкасты для учителей, набор демонстрационных программ разного уровня сложности. Технически возможным направлением является улучшение обратной связи при ошибках ассемблирования (указание строки и типовой причины), а также расширение диагностики при превышении лимита шагов, чтобы ученик понимал, что программа, вероятно, зациклилась. В целом проект оставляет пространство для эволюции без полной переделки архитектуры, что свидетельствует о разумности исходного проектного решения.

\newpage
\section{Источники}

\begin{enumerate}
  \item Федеральный государственный образовательный стандарт среднего общего образования (ФГОС СОО).
  \item Python Software Foundation. Python 3 Documentation. URL: \url{https://docs.python.org/3/} (дата обращения: 10.04.2026).
  \item React Team. React Documentation. URL: \url{https://react.dev/} (дата обращения: 10.04.2026).
  \item Microsoft. TypeScript Documentation. URL: \url{https://www.typescriptlang.org/docs/} (дата обращения: 10.04.2026).
  \item Vite Team. Vite Guide. URL: \url{https://vite.dev/guide/} (дата обращения: 10.04.2026).
  \item MOS Technology. \textit{MCS6500 Microcomputer Family Programming Manual}. 1976.
  \item Zaks R. \textit{Programming the 6502}. Sybex, 1981.
  \item Masswerk. 6502 Instruction Set Reference. URL: \url{https://www.masswerk.at/6502/} (дата обращения: 10.04.2026).
  \item MDN Web Docs. Fetch API. URL: \url{https://developer.mozilla.org/docs/Web/API/Fetch_API} (дата обращения: 10.04.2026).
  \item Caddy Server Documentation. URL: \url{https://caddyserver.com/docs/} (дата обращения: 10.04.2026).
  \item Git Documentation. URL: \url{https://git-scm.com/doc} (дата обращения: 10.04.2026).
\end{enumerate}

\end{document}
